% Evans second edition, section 2.2.
% English source: printed p. 19; PDF page 36 (1-based).
% Chinese section starts checked from chapter opening, PDF 31; mean-value pages checked at PDF 36-37 and 53-55.
% Status: editorial study draft with reconstructed proofs and mathematical annotations.
\section{Laplace's Equation}
\label{sec:2-2}

\subsection{Fundamental solution (Evans \S2.2.1, Theorem 1)}
A radial function $v(x)=F(r)$, $r=|x|>0$, satisfies
\[
 \Delta v=F''(r)+\frac{n-1}rF'(r).
\]
Thus harmonicity off zero gives $(r^{n-1}F')'=0$. Integration yields
a constant plus $r^{2-n}$ if $n\geq3$, or a constant plus $\log r$
if $n=2$. Choose the constants so that the inward flux is one:
\[
 \Phi(x)=
 \begin{cases}
 [n(n-2)\omega_n]^{-1}|x|^{2-n},&n\geq3,\\
 -(2\pi)^{-1}\log|x|,&n=2.
 \end{cases}
 \qquad -\Delta\Phi=\delta_0.
\]
\Needspace{45mm}
\PDEAnnotation{符号由通量固定}{我们求解 $-\Delta u=f$，故要求 $-\int_{\partial B_r}\partial_\nu\Phi=1$。若改用 $\Delta$ 作基本算子，基本解及 Green 公式的符号也要一起改。}
\begin{expandedproof}[the distributional singularity]
The displayed function is locally integrable and harmonic away from zero,
and $\partial_r\Phi=-1/(n\omega_n r^{n-1})$.
For $\phi\in C_c^\infty$, apply Green's identity on a large ball
containing its support, with $\overline B_\varepsilon$ removed:
\[
 \int_{|x|>\varepsilon}\Phi(-\Delta\phi)
 =\int_{\partial B_\varepsilon}
      \bigl(\Phi\partial_r\phi-\phi\partial_r\Phi\bigr)\,dS.
\]
The first term tends to zero: it is $O(\varepsilon)$ for $n\geq3$,
and $O(\varepsilon|\log\varepsilon|)$ for $n=2$.
The second is the normalized spherical mean of $\phi$, tending to
$\phi(0)$. Local integrability permits passage in the volume integral.
Hence $-\Delta\Phi=\delta_0$ as a distribution.
\end{expandedproof}

For $f\in C_c^2(\R^n)$ put $u=\Phi*f$. Write
$u(x)=\int\Phi(y)f(x-y)\,dy$ and differentiate $f$, not the singular
kernel. On each compact set of $x$, all these integrals have common
bounded support in $y$ and are dominated by an integrable multiple of
$|\Phi|$. Thus $u\in C^2$ and
\[
 -\Delta u(x)=\int\Phi(y)(-\Delta_y f(x-y))\,dy=f(x)
\]
by the distribution identity just proved. This avoids an illegal
unqualified differentiation of $D^2\Phi$ at zero.
For $n=1$, $\Phi(x)=-|x|/2$ has the same normalization.

\subsection{Mean-value formulas (Evans \S2.2.2, Theorems 2--3)}
\label{sec:2-elliptic-mean}
Write $\omega_n=|B_1|$ and $\sigma_{n-1}=|\partial B_1|=n\omega_n$.
For $\overline B_r(x)\subset U$ define the \emph{normalized} means
\[
 S_u(x,r)=\frac1{n\omega_n r^{n-1}}\int_{\partial B_r(x)}u\,dS,
 \qquad
 A_u(x,r)=\frac1{\omega_n r^n}\int_{B_r(x)}u\,dy.
\]
The sphere is a boundary of dimension $n-1$; the ball has dimension $n$.
These are two different averaging measures.
\Needspace{45mm}
\PDEAnnotation{平均不是积分本身}{球面平均除以面积，球体平均除以体积。对常数 $u=1$，两者都必须等于 $1$；漏掉 $r^{n-1}$ 或 $r^n$ 就会破坏这一检验。}

\begin{theorem}[Mean-value identities]\label{thm:2-harmonic-mean}
If $u\in C^2(U)$ and $\Delta u=0$, then
\[
 u(x)=S_u(x,r)=A_u(x,r)
 \quad\text{whenever }\overline B_r(x)\subset U.
\]
More generally, for any $C^2$ function,
\[
 \partial_r S_u(x,r)
 =\frac1{n\omega_n r^{n-1}}\int_{B_r(x)}\Delta u(y)\,dy .
 \tag{M}
\]
\end{theorem}
\begin{expandedproof}[all radius and surface factors]
\paragraph{1. Remove the moving surface before differentiating.}
Set $y=x+r\theta$ on the unit sphere. Its surface Jacobian is $r^{n-1}$,
so
\[
 S_u(x,r)=\frac1{n\omega_n}\int_{\partial B_1}u(x+r\theta)\,dS_\theta.
\]
On a fixed compact range of radii, $Du$ is continuous and bounded.
Differentiation under this fixed finite integral is legitimate:
\[
 \partial_r S_u(x,r)
 =\frac1{n\omega_n}\int_{\partial B_1}Du(x+r\theta)\cdot\theta\,dS_\theta.
\]
Transforming back, $\theta=\nu(y)$ is the outward unit normal, giving
\[
 \partial_r S_u(x,r)
 =\frac1{n\omega_n r^{n-1}}
       \int_{\partial B_r(x)}Du\cdot\nu\,dS.
\]
The divergence theorem on the smooth ball changes the flux into
$\int_{B_r(x)}\Delta u$. This proves (M).

\paragraph{2. Identify the constant.}
If $\Delta u=0$, the derivative in (M) vanishes for every admissible $r$.
Hence the spherical mean is independent of radius.
Continuity gives
\[
 |S_u(x,r)-u(x)|
 \leq\sup_{|y-x|\leq r}|u(y)-u(x)|\longrightarrow0.
\]
The constant is therefore exactly $u(x)$, not an unspecified constant.


\PDEMarginDiagram{2.M}{球体由球壳相加}
{先对每个半径做球面平均，再乘球壳体积并积分。体积平均的径向权重不是均匀的。}
{ch02-shell-margin}

\paragraph{3. Recover the ball average by adding spherical shells.}
Polar integration yields
\[
 \int_{B_r(x)}u
 =\int_0^r\int_{\partial B_\rho(x)}u\,dS\,d\rho
 =\int_0^r n\omega_n\rho^{n-1}u(x)\,d\rho
 =\omega_n r^n u(x).
\]
Divide by the volume. This also proves, for general continuous $u$,
\[
 A_u(x,r)=\frac n{r^n}\int_0^r \rho^{n-1}S_u(x,\rho)\,d\rho.
 \tag{SB}
\]
\end{expandedproof}
\Needspace{45mm}
\PDEAnnotation{方程如何变成平均}{关键中介是通量：$\Delta u=0$ 令每个球的总法向通量为零，通量又恰好是球面平均对半径的导数。故平均不随尺度变化。}

\paragraph{Can one recover spheres from balls?}
Yes. Differentiating $r^nA_u(x,r)$ in (SB) gives
\[
 S_u(x,r)=A_u(x,r)+\frac r n\partial_r A_u(x,r).
\]
Consequently a continuous function with the ball-mean identity for every
small radius also has the sphere-mean identity. A single radius at a
single center is not sufficient.

\begin{theorem}[Converse, including a continuous version]
If $u\in C^2(U)$ has either mean identity for every sufficiently small
ball at each center, then $\Delta u=0$.
If $u\in C(U)$ has this property at every center and all locally admissible
small radii, it is smooth and harmonic.
\end{theorem}
\begin{expandedproof}
For $C^2$ functions Taylor-expand at $x$ with a uniform $o(r^2)$ error.
Symmetry of the sphere cancels linear and mixed terms; its normalized
second moments are
$\mathbb E(\theta_i\theta_j)=\delta_{ij}/n$.
The ball has moments $\mathbb E(z_i z_j)=r^2\delta_{ij}/(n+2)$.
Thus
\[
 S_u(x,r)=u(x)+\frac{r^2}{2n}\Delta u(x)+o(r^2),
\]
\[
 A_u(x,r)=u(x)+\frac{r^2}{2(n+2)}\Delta u(x)+o(r^2).
 \tag{T}
\]
An exact identity for arbitrarily small $r$ forces $\Delta u(x)=0$.

For a merely continuous $u$, first turn ball means into sphere means
using (SB), if necessary. Take a radial smooth mollifier of mass one,
supported in $B_\varepsilon$. In polar coordinates,
\[
 (\rho_\varepsilon*u)(x)
 =\int_0^\varepsilon \rho_\varepsilon(r)
       \int_{\partial B_r(x)}u\,dS\,dr=u(x).
\]
The convolution is smooth on the smaller domain; hence $u$ is smooth
there and throughout $U$. Now apply the $C^2$ converse.
This proves the smoothness assertion without differentiating an
initially nonsmooth function.
\end{expandedproof}
\Needspace{45mm}
\PDEAnnotation{二阶误差读出 Laplace}{普通函数的平均一般不等于中心值；首个对称性允许的误差就是 $r^2\Delta u$。调和性让这一误差及全部更高阶平均误差精确抵消。}

\paragraph{Subharmonic and superharmonic functions.}
If $\Delta v\geq0$, (M) implies $S_v(x,r)\geq v(x)$ and then
$A_v(x,r)\geq v(x)$. If $\Delta v\leq0$, reverse both inequalities.
For $-\Delta u=f$, integrating (M) and changing the order of integration
gives, for $n\geq3$,
\[
 u(x)=S_u(x,r)+\int_{B_r(x)}
 \frac{|y-x|^{2-n}-r^{2-n}}{n(n-2)\omega_n}f(y)\,dy.
 \tag{P}
\]
Indeed $\int_{|y-x|}^r\rho^{1-n}\,d\rho
=(|y-x|^{2-n}-r^{2-n})/(n-2)$.
For $n=2$ replace the kernel by $(2\pi)^{-1}\log(r/|y-x|)$.
The kernel is nonnegative: a positive source raises $u(x)$ above
its boundary average under the convention $-\Delta u=f$.

\paragraph{Concrete checks.}
In two dimensions, $u(y)=y_1^2-y_2^2$ and $u(y)=y_1y_2$ have
$\Delta u=0$, so both means reproduce $u(x)$ at every center.
For $v(y)=|y|^2$,
\[
 S_v(0,r)=r^2,\qquad A_v(0,r)=\frac n{n+2}r^2,\qquad v(0)=0.
\]
For $n=1$, the sphere consists of two endpoints:
$S_u(x,r)=[u(x-r)+u(x+r)]/2$. Harmonic functions are affine,
and the same identities remain valid with $\omega_1=2$.
\Needspace{45mm}
\PDEAnnotation{别混同一元中值定理}{这里是积分平均精确重现中心值，不是声称“存在某个中间点”。后面的热球公式也保持中心值，但必须改变平均的几何区域和权重。}


\subsection{Consequences of averaging (Evans \S2.2.3)}
\begin{theorem}[Strong maximum principle; source Theorems 4--5]
If $U$ is bounded and $u\in C^2(U)\cap C(\overline U)$ is harmonic,
its maximum and minimum are attained on $\partial U$.
If $U$ is connected and a maximum is attained inside, $u$ is constant.
Consequently the classical Dirichlet problem $-\Delta u=f$, $u=g$
has at most one solution in this class.
\end{theorem}
\begin{expandedproof}
Let $M=\max_{\overline U}u$ and suppose $u(x)=M$ at an interior point.
The ball identity gives $\int_{B_r(x)}(M-u)=0$.
The integrand is nonnegative and continuous. If it were positive at
one point, continuity would give a positive integral on a smaller ball,
a contradiction. Thus $u=M$ throughout $B_r(x)$.
The set $\{u=M\}$ is open and relatively closed in $U$.
On a connected component containing an interior maximum it is the
whole component, and continuity transfers this value to that component's
boundary, which lies in $\partial U$. This proves the boundary maximum
statement even if $U$ has several components.
Apply the argument to $-u$ for minima.
The difference of two Dirichlet solutions is harmonic with zero boundary
values, so its maximum and minimum are both zero.
\end{expandedproof}
\Needspace{45mm}
\PDEAnnotation{正权平均的等号条件}{平均值达到所有被平均数的上界，意味着它们几乎处处都达到该上界；连续性再把“几乎处处”提升为处处。热球的强最大值原理使用同一机制。}

\paragraph{Smoothness and derivative estimates (source Theorems 6--7).}
Radial convolution equals $u$ by the mean identity, as proved in
\S\ref{sec:2-elliptic-mean}, so harmonic functions are $C^\infty$.
Fix a radial mollifier $\rho$ with support in $B_1$.
For $B_r(x)\Subset U$, use a convolution radius $r/2$ near $x$:
\[
 |D^\alpha u(x)|
 =\left|\int D^\alpha\rho_{r/2}(x-y)u(y)\,dy\right|
 \leq C_{\alpha,n}r^{-n-|\alpha|}\|u\|_{L^1(B_r(x))}.
 \tag{D}
\]
All derivatives land on the smooth kernel; its scaling explains the power
of $r$. With a first derivative the same argument gives
$\|Du\|_{L^\infty(B_{R-\delta})}
 \leq C_n\delta^{-1}\|u\|_{L^\infty(B_R)}$.

\begin{theorem}[Liouville and bounded Poisson representation; source 8--9]
A bounded harmonic function on $\R^n$ is constant.
For $n\geq3$ and $f\in C_c^2$, every bounded solution of
$-\Delta u=f$ is $\Phi*f+C$.
\end{theorem}
\begin{expandedproof}
Apply (D) for $|\alpha|=1$ on $B_R(x)$:
$|Du(x)|\leq C_nR^{-1}\|u\|_\infty\to0$ as $R\to\infty$.
Thus $Du=0$ and connectedness makes $u$ constant.
For $n\geq3$, $\Phi*f$ is continuous and tends to zero at infinity,
hence bounded. Subtract it from $u$ and apply Liouville.
For $n=2$, the logarithmic potential need not be bounded at infinity,
so that final step cannot be asserted with the same hypotheses.
\end{expandedproof}

\begin{theorem}[Analyticity; source Theorem 10]
Harmonic functions are real analytic in the interior.
\end{theorem}
\begin{expandedproof}[quantitative iteration rather than smoothness alone]
Fix $B_R(x_0)\Subset U$, with $M=\|u\|_{L^\infty(B_R)}$.
Every derivative of $u$ is harmonic, because derivatives commute with
$\Delta$. Iterate the first-derivative estimate $k$ times on concentric
balls whose successive radii differ by $R/(2k)$.
For any $|\alpha|=k$ this gives
\[
 \|D^\alpha u\|_{L^\infty(B_{R/2})}
 \leq M(2C_n k/R)^k
 \leq M(2eC_n/R)^k k!.
\]
The last inequality follows from
$\log(k!)=\sum_{j=1}^k\log j\geq\int_1^k\log s\,ds
\geq k\log k-k$.
For $h$ small enough that the segment from $x_0$ to $x_0+h$ stays in
$B_{R/2}$, apply one-dimensional Taylor's theorem to
$q(t)=u(x_0+th)$. Its $k$th derivative is
\[
 q^{(k)}(t)=\sum_{|\alpha|=k}\frac{k!}{\alpha!}
                    h^\alpha D^\alpha u(x_0+th).
\]
Dividing the remainder by $k!$ and using
$\sum_{|\alpha|=k}|h^\alpha|/\alpha!=|h|_1^k/k!$ bounds it by
$M(2eC_n|h|_1/R)^k$, tending to zero for sufficiently small $h$.
Thus the multivariable Taylor series equals $u$ near $x_0$.
\end{expandedproof}
\Needspace{45mm}
\PDEAnnotation{光滑不等于解析}{$C^\infty$ 只说明每阶导数存在；解析还需控制高阶导数的增长。这里的 $k!$ 型估计让 Taylor 余项趋零，不能省略这一环。}

\begin{theorem}[Harnack; source Theorem 11]
For connected $V\Subset U$ and nonnegative harmonic $u$,
$\sup_Vu\leq C_{U,V,n}\inf_Vu$.
\end{theorem}
\begin{expandedproof}
Choose $\rho>0$ with every $B_{2\rho}(z)$, $z\in V$, inside $U$.
If $x,y\in V$ and $|x-y|\leq\rho$, then
$B_\rho(y)\subset B_{2\rho}(x)$ and positivity gives
\[
 u(x)=\frac1{|B_{2\rho}|}\int_{B_{2\rho}(x)}u
 \geq2^{-n}\frac1{|B_\rho|}\int_{B_\rho(y)}u=2^{-n}u(y).
\]
Exchange $x,y$ for the other inequality.
Compactness gives a finite $\rho/4$ net of $V$.
Connectedness permits a chain between any two net points with successive
centers separated by at most $\rho$; one may obtain this by covering a
path in $V$ and deleting loops in the finite cover graph.
There is a uniform finite bound on the chain length.
Multiplying the local factors proves the global comparison.
If one value is zero, the same inequalities force every value in $V$
to be zero, so the statement includes that case.
\end{expandedproof}

\subsection{Green and Poisson kernels (Evans \S2.2.4)}
For sufficiently smooth $u,v$, the two Green identities are
\[
 \int_U Du\cdot Dv=-\int_Uv\Delta u+\int_{\partial U}v\partial_\nu u,
\]
\[
 \int_U(u\Delta v-v\Delta u)
 =\int_{\partial U}(u\partial_\nu v-v\partial_\nu u).
\]
The second is the difference of the first identity in the two orders.

\paragraph{Boundary correction.}
For fixed $x\in U$, suppose a harmonic corrector $h_x$ with
$h_x(y)=\Phi(y-x)$ on $\partial U$ has been constructed.
Set $G(x,y)=\Phi(y-x)-h_x(y)$. Then
\[
 -\Delta_yG(x,y)=\delta_x,\qquad G(x,y)=0\quad(y\in\partial U).
\]
General-domain existence of this corrector is not proved by the notation;
it belongs to Dirichlet existence theory later in the book.
The following representation is conditional on such a sufficiently
regular Green function. Explicit examples below avoid that existence issue.
\Needspace{45mm}
\PDEAnnotation{Green 函数是修正后的点源}{基本解有正确的点源奇性，却通常不满足边界条件。减去调和修正项保留奇性、消掉边界值；“定义修正项”本身不等于证明它存在。}

\begin{theorem}[Representation and symmetry; source Theorems 12--13]
For $-\Delta u=f$, $u=g$ on the boundary,
\[
 u(x)=\int_U G(x,y)f(y)\,dy
       -\int_{\partial U}g(y)\partial_{\nu_y}G(x,y)\,dS_y.
 \tag{G}
\]
Furthermore $G(x,y)=G(y,x)$.
\end{theorem}
\begin{expandedproof}
Apply the second Green identity on $U\setminus\overline B_\varepsilon(x)$
to $u$ and $G(x,\cdot)$. Away from $x$, $\Delta G=0$.
On the outer boundary $G=0$, leaving $\int g\partial_\nu G$.
On the inner boundary the singular normal derivative has total flux
one, with normal directed into the removed ball; the resulting term
tends to $u(x)$. Terms involving $G\partial_\nu u$ tend to zero at
rates $O(\varepsilon)$ or $O(\varepsilon|\log\varepsilon|)$.
Thus $\int_U Gf=\int_{\partial U}g\partial_\nu G+u(x)$, giving (G).

For symmetry, apply the same second identity to
$G(x,\cdot),G(z,\cdot)$ with small disjoint balls about $x,z$ removed.
Both are harmonic on the remaining domain and zero on the outer
boundary. At the sphere around $x$, the term with the singular normal
derivative tends to $G(z,x)$ and the opposite product tends to zero;
the sphere around $z$ gives $-G(x,z)$.
Their sum is zero, proving symmetry.
\end{expandedproof}

\paragraph{The upper half-space (source Theorem 14).}
For $x_n>0$ let $x^*=(x',-x_n)$.
Reflection gives $G(x,y)=\Phi(y-x)-\Phi(y-x^*)$.
The singularity of the image lies outside the domain, and the two
distances agree on $y_n=0$.
Its boundary normal derivative gives
\[
 P(x,z)=\frac{2x_n}{n\omega_n(|x'-z|^2+x_n^2)^{n/2}},
 \quad z\in\R^{n-1}.
\]
\[
 u(x)=\int_{\R^{n-1}}P(x,z)g(z)\,dz.
\]
The kernel is harmonic in $x$, positive, and has mass one.
To check the mass without imposing boundary decay, put $z=x'+x_n w$;
the remaining integral is
$[2/(n\omega_n)]\int(1+|w|^2)^{-n/2}\,dw=1$.
Its value follows by polar coordinates, or the Gaussian integral and
Gamma substitution.
For bounded continuous $g$, split the boundary integral into
$|z-x'_0|<\delta$ and its complement.
Uniform continuity near $x'_0$ controls the first error; the latter
kernel mass tends to zero as $x\to(x'_0,0)$.
This proves boundary recovery and, by differentiation on compact interior
sets, harmonicity.
Uniqueness here is among bounded solutions: a bounded zero-boundary
harmonic function extends oddly across the plane, is harmonic by the
reflection argument in Problem 10, and is constant by Liouville, hence zero.


\PDEBodyDiagram{2.2}{A reflected point source enforces the boundary value.}
{The solid source at $x$ is inside the upper half-space; the image at $x^*$ is outside and has the opposite sign. Every boundary point $y$ is equidistant from the pair, so their potentials cancel there. The labels indicate source signs, not signed coordinate locations.}
{ch02-image-source}

\paragraph{The ball (source Theorem 15).}
For $n\geq3$, $0<|x|<R$, put $x^*=R^2x/|x|^2$. An explicit Green function is
\[
 G(x,y)=\Phi(y-x)-(R/|x|)^{n-2}\Phi(y-x^*).
\]
On $|y|=R$, $|y-x^*|=(R/|x|)|y-x|$, so the boundary value is zero.
At $x=0$ the formula has the limit $\Phi(y)-\Phi(R)$.
In two dimensions use instead
\[
 G(x,y)=\frac1{2\pi}
 \log\frac{|x|\,|y-x^*|}{R|y-x|},
\]
with the corresponding center limit $(2\pi)^{-1}\log(R/|y|)$.
Differentiating in the boundary normal gives, for every $n\geq2$,
\[
 P_R(x,y)=\frac{R^2-|x|^2}{n\omega_n R|x-y|^n},
 \qquad
 u(x)=\int_{\partial B_R}P_R(x,y)g(y)\,dS_y.
 \tag{PB}
\]
\begin{expandedproof}[why the boundary formula solves the problem]
The constructed Green function and (G) applied to $u=1$ give
$\int_{\partial B_R}P_R=1$. This uses a known constant harmonic function,
so does not assume existence for arbitrary data.
For $|y|=R$, direct differentiation yields
\[
 \Delta_x\big[(R^2-|x|^2)|x-y|^{-n}\big]
 =2n(R^2-|y|^2)|x-y|^{-n-2}=0.
\]
Thus the integral is harmonic; all differentiations are dominated
on compact subsets separated from the sphere.
For $x\to y_0\in\partial B_R$ split $g(y)-g(y_0)$ into a small spherical
neighborhood and its complement. Positivity and mass one control the
near part by the modulus of continuity. On the complement the denominator
stays bounded below while $R^2-|x|^2\to0$, so the far part tends to zero.
This proves the boundary data. The maximum principle gives uniqueness.
At $x=0$, $P_R$ is constant and (PB) reduces exactly to the spherical
mean identity.
\end{expandedproof}
\Needspace{45mm}
\PDEAnnotation{中心平均与偏心平均}{从球心看，边界各方向权重相同；从偏心点看，靠近该点的边界权重更大，变成 Poisson 核。球面均值是 Poisson 表示的球心特例。}

\subsection{Energy and Dirichlet's principle (Evans \S2.2.5, Theorems 16--17)}
Let $U$ be bounded and smooth enough for the following integrations.
For two $C^2(\overline U)$ solutions of the same Dirichlet problem,
their difference $w$ has $\Delta w=0$ and $w|_{\partial U}=0$.
Green's first identity gives
$\int_U|Dw|^2=-\int_Uw\Delta w+\int_{\partial U}w\partial_\nu w=0$.
Thus $w$ is constant on each component, and its boundary values make it zero.

For $-\Delta u=f$, $u=g$, define
\[
 J(v)=\int_U\left(\tfrac12|Dv|^2-fv\right)\,dx.
\]
For $v=u+w$ with $w|_{\partial U}=0$,
\[
 J(v)-J(u)
 =\int_U(Du\cdot Dw-fw)+\tfrac12\int_U|Dw|^2
 =\tfrac12\int_U|Dw|^2\geq0.
\]
Here the linear term vanishes by Green's identity and the PDE.
Conversely, if a smooth $u$ minimizes $J$ for its boundary data, test
$v=u+\varepsilon\phi$ with $\phi\in C_c^\infty(U)$ and both signs of
$\varepsilon$. Differentiating the quadratic polynomial at zero gives
$\int Du\cdot D\phi=\int f\phi$, hence $-\Delta u=f$.
This proves the equivalence in the smooth class; existence of a minimizer
in a Sobolev class uses the later functional-analytic theory.
\Needspace{45mm}
\PDEAnnotation{从第二章走向第五章}{能量只含一阶导数，所以弱形式 $\int Du\cdot D\phi=\int f\phi$ 自然进入 $H^1$。第二章的显式公式与第五章的函数空间方法在这里接上。}
